\typeout{NT FILE chapter10.tex}%

\chapter{Conclusion and Future Work}
\label{cha:conclusion}

\section{Conclusion}
\label{sec:conclusion_summary}

This dissertation extended both the new improved OCamlFLAT library and the OFLAT web application with full support for finite-state
transducers (FST). On the library side, we added output-aware acceptance functions (\texttt{acceptOut},
\texttt{acceptExpectedOutput}), minimization and determinism checks adapted to input/output labels, and utilities used by
the UI (e.g., \texttt{inputEdges}). On the UI side, OFLAT now edits and executes FSTs with the same interaction patterns as
other models: input:output edge labels, step-by-step execution with branch highlights, accumulated-output display in the
info box, popper counts for nondeterminism, synchronized trace/transition tables, and exercise validation using the
\texttt{input->output} format. The bridge built with \texttt{js\_of\_ocaml} keeps the browser client in sync with the
OCaml models.

Evaluation combined unit tests in OCamlFLAT with manual and integration checks in OFLAT, exercising deterministic and
nondeterministic runs, Moore-style outputs, and minimization. The result is a cohesive teaching tool where students can
construct, visualize, step through, and assess FSTs with immediate feedback that matches formal definitions.

\section{Summary of Contributions}
\label{sec:conclusion_contrib}
\begin{itemize}
	\item \textbf{Library}: New FST data model and algorithms 
    (acceptOut/acceptExpectedOutput, determinization variants, minimization, 
    composition/union/intersection/inverse/concatenation, projection, TM conversion) 
    plus JSON interoperability for persistence.
	\item \textbf{UI}: End-to-end FST authoring and execution in OFLAT with input:output edge labels, 
    trail-driven stepping, output-aware exercises, and reuse of Cytoscape-based visualization primitives.
	\item \textbf{Pedagogy}: Consistent interaction patterns across FA/PDA/TM/FST, 
    output-aware checking for exercises, and tables/traces aligned with textbook definitions.
	\item \textbf{Tooling}: Single-language pipeline via \texttt{js\_of\_ocaml}, 
    keeping logic and UI behavior synchronized.
\end{itemize}

\section{Future Work}
\label{sec:future_work}

The work conducted in this dissertation will be integrated into the latest OFLAT release in the near future,
this new version has various improvements and bug fixes on the UI side, but the modular approach of this implementation
allows this integration to be seamless.

Besides that, several extensions could further improve the platform:
\begin{itemize}
    \item \textbf{Exercise and grading flow:} support automatic generation of input/output test sets and richer feedback for
        counterexamples.
    \item \textbf{UX polish and accessibility:} Colorblind-safe palettes, and add
        lightweight tutorials or guided tours for first-time users.
    \item \textbf{Further models:} explore other transducer models (e.g., weighted, probabilistic) to expand the didactic scope while reusing
        the existing interaction patterns.
    \item \textbf{Performance and scale:} stress-test high-branching FSTs, add non-termination guards in determinization, 
	and show clear guidance on practical limits.
\end{itemize}
